problem:  
Let \( V^n \subset S^{n+k} \) be a stationary integral \( n \)-varifold (so with zero generalized mean curvature) whose regular part is a smooth minimal submanifold and whose second fundamental form \( A \) satisfies \( |A| \in L^2(\mathrm{reg}\, V) \).  
Assume that every tangent cone to \( V \) at a singular point is a minimizing cone in \( \mathbb{R}^{n+k+1} \), and that each such cone’s link \( L \subset S^{n+k} \) is a minimal submanifold of the round sphere.  

**Question:**  
Does the positive curvature of the ambient sphere impose a quantitative restriction on the possible links \( L \)?  
More precisely, is it true that for each fixed \( n \), there exists a constant \( \lambda_n > 0 \) such that if \( L \subset S^{n+k} \) arises as the link of a tangent cone to any stationary integral \( n \)-varifold in \( S^{n+k} \) satisfying the \( L^2 \) curvature bound, then the first nonzero eigenvalue of the Laplace–Beltrami operator on \( L \) satisfies  
\[
\lambda_1(L) \ge \lambda_n?
\]  
If such a universal spectral gap existed, it would force the cone to be linearly stable only for a restricted family of links—potentially implying that all tangent cones must be spherical and thus that \(V\) is smooth.  
If no such universal bound exists, can one construct an explicit sequence of minimal varifolds in round spheres with increasingly degenerate links exhibiting collapsing eigenvalues?

Why is it a "good" problem:  
This problem refines the vague hope that positive ambient curvature might restrict singular minimal configurations into a quantitative geometric form: an a priori spectral gap on cone links. It is logically well-posed within geometric measure theory—asking about the spectral geometry of links of tangent cones under \(L^2\)-bounded curvature hypotheses, a meaningful regime in which existing tools of Almgren-type regularity and Simons’ identities can interact fruitfully with eigenvalue estimates.  

It bridges *regularity theory of minimal varifolds* (microlocal analysis of tangent cones) with *global spectral geometry* (eigenvalues on minimal submanifolds of spheres).  
If true, it would uncover a new rigidity phenomenon for singularities in positively curved spaces; if false, the construction of counterexamples would illuminate the analytic spectrum of singular minimal submanifolds.  

The statement is simple and geometric—phrased entirely in terms of minimal varifolds, tangent cones, and Laplacian eigenvalues—yet connects deep analytic, geometric, and variational structures, embodying the “narrow entrance, profound depth” quality of a good research problem.